% ==================================================================
% YangTex · 编译产物，请勿直接编辑（改动会在下次保存时被覆盖）
% 由 Markdown 源 main.md 自动生成
% 请编辑同目录下的 main.md / main.mdtex：那才是主文档，
% 用的是 Markdown 语法；本文件只交给 TeX 引擎编译。
% ==================================================================
\documentclass[UTF8,a4paper,11pt,fontset=fandol]{ctexart}
\usepackage{amsmath}
\usepackage{amssymb}
\usepackage{mathtools}
\usepackage{amsthm}
\usepackage{graphicx}
\usepackage{xcolor}
\usepackage{enumitem}
\usepackage{geometry}
\usepackage{hyperref}
\geometry{margin=2.5cm}
\newtheorem{theorem}{定理}[section]
\newtheorem{lemma}[theorem]{引理}
\newtheorem{definition}[theorem]{定义}
\newtheorem{proposition}[theorem]{命题}
\newtheorem{corollary}[theorem]{推论}
\title{未命名文档}
\author{}
\date{\today}
\begin{document}
\maketitle

\% ================================================================== \% YangTex · 由 Markdown 模板自动生成（内置 Markdown -\textgreater{} LaTeX 转换器） \% 源模板：md\_meeting\_minutes.md \% 这是标准 LaTeX：可直接编辑，实时编译 / SyncTeX 双向同步照常可用。 \% 需要继续用 Markdown 写时，选中文本后点「MD→TeX 转换」。 \% ================================================================== \textbackslash{}documentclass[UTF8,a4paper,11pt,fontset=fandol]\{ctexart\} \textbackslash{}usepackage\{amsmath\} \textbackslash{}usepackage\{amssymb\} \textbackslash{}usepackage\{mathtools\} \textbackslash{}usepackage\{amsthm\} \textbackslash{}usepackage\{graphicx\} \textbackslash{}usepackage\{xcolor\} \textbackslash{}usepackage\{enumitem\} \textbackslash{}usepackage\{geometry\} \textbackslash{}usepackage\{hyperref\} \textbackslash{}geometry\{margin=2.5cm\} \textbackslash{}newtheorem\{theorem\}\{定理\}[section] \textbackslash{}newtheorem\{lemma\}[theorem]\{引理\} \textbackslash{}newtheorem\{definition\}[theorem]\{定义\} \textbackslash{}newtheorem\{proposition\}[theorem]\{命题\} \textbackslash{}newtheorem\{corollary\}[theorem]\{推论\} \textbackslash{}title\{组会纪要\} \textbackslash{}author\{数学与统计学院\} \textbackslash{}date\{\textbackslash{}today\} \textbackslash{}begin\{document\} \textbackslash{}maketitle \textbackslash{}section\{组会基本信息\}

\textbackslash{}begin\{itemize\} \textbackslash{}item \textbackslash{}textbf\{时间\}：每周三 15:00 \textbackslash{}item \textbackslash{}textbf\{地点\}：理科楼研讨室 \textbackslash{}item \textbackslash{}textbf\{主题\}：不可压缩流体方程的正则性 \textbackslash{}end\{itemize\}

\textbackslash{}subsection\{本次汇报\}

\textbackslash{}begin\{enumerate\} \textbackslash{}item 博士生甲：能量估计与紧性论证进展 \textbackslash{}item 博士生乙：数值格式稳定性测试 \textbackslash{}item 教师小结：下一步重点 \textbackslash{}end\{enumerate\}

\textbackslash{}subsection\{问题与结论\}

\textbackslash{}begin\{quote\} 先验估计中的常数 $C$ 需统一口径，避免与 $\nu$ 混淆。 \textbackslash{}end\{quote\}

\textbackslash{}begin\{center\} \textbackslash{}begin\{tabular\}\{l\textbar{}l\textbar{}l\textbar{}l\} \textbackslash{}hline 事项 \& 负责人 \& 截止 \& 状态 \textbackslash{}\textbackslash{} \textbackslash{}hline 补齐 $n=3$ 的估计 \& 甲 \& 本周五 \& 进行中 \textbackslash{}\textbackslash{} 复现算例 2 \& 乙 \& 下周 \& 未开始 \textbackslash{}\textbackslash{} 整理参考文献 \& 全体 \& 月底 \& 进行中 \textbackslash{}\textbackslash{} \textbackslash{}hline \textbackslash{}end\{tabular\} \textbackslash{}end\{center\}

\textbackslash{}subsection\{下次议程\}

\textbackslash{}begin\{itemize\} \textbackslash{}item 检查 $\|\nabla u\|_{L^2}$ 的一致性估计 \textbackslash{}item 讨论投稿期刊选择 \textbackslash{}item 确定课程讲义分工 \textbackslash{}end\{itemize\}

\textbackslash{}par\textbackslash{}noindent\textbackslash{}rule\{\textbackslash{}textwidth\}\{0.4pt\}\textbackslash{}par

纪要由 YangTex Markdown 模板生成，表格与列表可直接照常编辑。

\textbackslash{}end\{document\}

\end{document}
